\documentclass[12pt,a4paper]{article}
\usepackage[utf8]{inputenc}
\usepackage[T1]{fontenc}
\usepackage[margin=2.5cm]{geometry}
\usepackage{graphicx}
\usepackage{listings}
\usepackage{xcolor}
\usepackage{caption}

\newcommand{\studentgroup}{FAF-242}
\definecolor{termbg}{rgb}{0.96,0.96,0.96}
\lstset{backgroundcolor=\color{termbg},basicstyle=\ttfamily\footnotesize,breaklines=true,frame=single}

\begin{document}
\begin{titlepage}
    \centering
    \IfFileExists{utm_logo.png}{\includegraphics[width=0.25\textwidth]{utm_logo.png}\\[0.6cm]}{\vspace*{2cm}}
    {\large \textbf{UNIVERSITATEA TEHNICĂ A MOLDOVEI}}\\[0.2cm]
    {\large Facultatea Calculatoare, Informatică și Microelectronică}\\[0.2cm]
    {\large Departamentul Ingineria Software și Automatică}\\[2.5cm]
    {\Large \textbf{LUCRAREA DE LABORATOR NR. 8}}\\[0.4cm]
    {\LARGE \textbf{Aceleași apeluri, în OS-ul tău (xv6)}}\\[0.4cm]
    {\large Disciplina: Sisteme de Operare}\\[2.5cm]
    \begin{minipage}{0.45\textwidth}\flushleft\textbf{Done by:}\\st. gr. \studentgroup\\Daniel-Ioan \textsc{Strunga}\end{minipage}\hfill
    \begin{minipage}{0.50\textwidth}\flushright\textbf{Checked by:}\\asist. univ. Patricia \textsc{Reițman}\\asist. univ. Marius \textsc{Dumitrașcu}\end{minipage}
    \vfill{\large Chișinău --- 2026}
\end{titlepage}
\newpage

\section{Scopul lucrării}
Implementarea administrării de procese interdependente (IPC) și realizarea de comunicații secvențiale bidirecționale în interiorul sistemului operațional educațional xv6, verificând transpunerea conceptelor arhitecturii POSIX studiate la Laboratorul 7.

\section{Desfășurarea lucrării}
\begin{figure}[h!]
\begin{lstlisting}[language=bash]
$ pingpong
10: received ping
9: received pong
\end{lstlisting}
\caption{Figure 8.1 -- Comunicarea bidirecțională sincronizată între procese în xv6}
\end{figure}

\subsection{Răspunsuri (Observe)}
\begin{enumerate}
    \item \textbf{Part 1 (\texttt{run}):} Apelul sistemului subiacent a rămas neschimbat (triada \texttt{fork -> exec -> wait}). Procesul părinte își duplică zona de memorie, copilul rezultat își schimbă ramura curentă cu apelul \texttt{exec}, chemând intern utilitarul \texttt{echo}, care devine responsabil pentru printarea mesajului pe ecran.
    \item \textbf{Part 2 (\texttt{pingpong}):} Pipe-urile în mediile derivate din Unix (inclusiv xv6) sunt structuri de memorie \textit{unidirecționale}. Prin urmare, fluxul bidirecțional ("ping" urmat de un răspuns "pong") necesită crearea simultană a două instanțe separate de \texttt{pipe}, una alocată exclusiv pentru scrierea copilului, și alta alocată exclusiv pentru citirea acestuia.
\end{enumerate}

\section{Concluzii (Comparație Linux vs xv6)}
În realizarea lucrărilor 7 și 8 am observat universalitatea apelurilor POSIX. Principiul de separare logică \texttt{fork} + \texttt{exec} a rămas neatins între Linux (arhitectură de producție) și xv6 (kernel educațional minimal). 

Diferențele apar doar în complexitatea bibliotecilor care interfațează kernelul: Linux (prin glibc) adaugă opțiuni avansate (ex: \texttt{execlp} pentru căutarea automată în calea curentă sau biții de status suplimentari din funcția \texttt{waitpid} necesari decodificării întreruperilor hardware), în timp ce xv6 oferă strict rutinele esențiale neprocesate, forțând programatorul să rezolve rutarea și sincronizările inter-procese la nivel absolut elementar.
\end{document}
